\section{Redundancy five: exceptional cubic covers}\label{sec:r5}

Here $f=(a_0:\cdots:a_4)$ and $W_f$ is the pencil annihilated by
\[
 \begin{pmatrix}a_0&a_1&a_2&a_3\\a_1&a_2&a_3&a_4\end{pmatrix}.
\]
These are divided-power coordinates; the coefficients of a cubic in
$W_f$ are ordinary symmetric-power coordinates.  Factoring the formal
quartic $\sum a_iT^iU^{4-i}$ would therefore be noninvariant in
characteristics two and three.

\paragraph{The criterion in finite-geometric form.}
At $r=5$ the witness system $W_f\subset\Sym^3E^\vee$ is a pencil of
binary cubics, that is a line of $\mathrm{PG}(3,q)$ relative to the twisted
cubic; and a completely split squarefree cubic is a point of the
$\PGL_2(q)$-orbit $\mathcal O_3=\PGL_2(q)\cdot XY(X-Y)$ of size
$(q^3-q)/6$, equivalently a plane meeting the twisted cubic in three
distinct rational points \cite[Prop.~3.1 and Rem.~3.2]{BPS2022},
\cite[Lem.~4.1]{KaipaPradhan2025}.  Corollary~\ref{cor:splitfree}
therefore reads
\begin{equation}\label{eq:r5-orbit-criterion}
       f\ \text{split-free}
       \iff
       W_f\cap\mathcal O_3=\varnothing .
\end{equation}
In this form the redundancy-five question is a point-orbit
distribution question for line orbits of $\mathrm{PG}(3,q)$, and as such it has
been answered in the finite-geometry literature; the relation is set
out in Remark~\ref{rem:r5-line-literature}.  What that literature does
not supply, and what this section adds, is the passage from a
redundancy-five syndrome to its pencil: which lines arise as some
$W_f$, and which \emph{syndromes} are thereby split-free.

\begin{proposition}[Radius gate]\label{prop:r5-radius}
For every prime power $q\geq7$,
\[
       \rho(\PRS(q-4))=4.
\]
\end{proposition}

\begin{proof}
The coding/arc criterion identifies radius four with completeness of
the quartic normal rational curve. The Seroussi--Roth extension theorem
\cite[Thm.~1]{SeroussiRoth1986}, applied to the dual five-dimensional
GRS code, gives nonextendability when
\[
       5\leq\left\lfloor\frac q2\right\rfloor+2,
\]
apart from its even-field dimension-three exception. The inequality
holds for \(q\geq7\), and the exception does not occur. Dür's
completeness--covering-radius equivalence \cite{Dur1994}, as stated by
Kaipa \cite[Sec.~IV and Prop.~4(1)]{Kaipa2017}, now gives radius four. Thus
Corollary~\ref{cor:splitfree} applies.  The finite certificate
independently checks the same conclusion only in its declared finite
domain; it is not the all-field radius proof.
\end{proof}

\subsection{Gcd strata}

\begin{proposition}[Quadratic gcd]\label{prop:r5-gcd2}
If $\deg\gcd(W_f)=2$, then
\[
       W_f=Q\langle T,U\rangle.
\]
If $Q$ is split squarefree, $f$ is shallow.  If $Q$ is irreducible,
$f$ belongs to the conjugate-secant family; if $Q=\lambda^2$, it
belongs to the tangent family.  Conversely every point in those
families has the displayed gcd.  Their deep-family sizes are
\[
       |\mathrm T|=q(q+1),\qquad
       |\mathrm S|=\frac{q(q^2-1)}2.
\]
\end{proposition}

\begin{proof}
The equality follows by dimension.  A split $Q$ and a linear cofactor
avoiding its two roots give a split squarefree cubic.  An irreducible
$Q$ survives in every member, while a square $Q$ forces a repeated
root in every member.  Lemma~\ref{lem:hankel}, including its confluent
and conjugate-root forms, identifies the corresponding syndromes with
the secant and tangent spans.  The sizes are the established
tangent/conjugate-secant counts from the lower-redundancy
classification
\cite[Prop.~1]{Kaipa2017} and \cite[Thms.~I.4--I.6]{ZWK2020}.
\end{proof}

\begin{proposition}[Linear gcd]\label{prop:r5-gcd1}
Suppose $\deg\gcd(W_f)=1$.  Then $f$ is shallow for $q\geq8$.  At
$q=7$ the same conclusion is part of Certificate R5.
\end{proposition}

\begin{proof}
Write $W_f=\lambda_0V$ with $V$ a basepoint-free quadratic pencil.  If
its degree-two map $\psi:\PP^1\to\PP^1$ is separable, its rational deck
involution $\iota$ gives the graph
\[
 \Gamma_\iota(\F_q)=\{(t,\iota(t)):t\in\PP^1(\F_q)\},
\]
which has exactly $q+1$ rational points.  A safe union bound deletes at
most two fixed graph points, at most four graph points over the branch
fibers, and at most the two graph points
$(\lambda_0,\iota(\lambda_0))$ and
$(\iota(\lambda_0),\lambda_0)$.  Thus fewer than $q+1$ graph points
are deleted when $q\geq8$.  A surviving point gives a split quadratic
with two distinct roots avoiding $\lambda_0$, and hence a split
squarefree cubic in $W_f$.  In characteristic two an inseparable
quadratic pencil is equivalent to $\langle T^2,U^2\rangle$; the two
overlapping Hankel equations then force $f$ onto the normal rational
curve, contrary to the rank-two pencil hypothesis.
\end{proof}

\subsection{The trivial-gcd cubic cover}

Choose a basis $c_1,c_2$ of a trivial-gcd pencil and put
\[
 X_f=\{([u:v],t):uc_1(t)+vc_2(t)=0\}
       \subset\PP^1\times\PP^1.
\]

\begin{proposition}[Cubic incidence and off-diagonal fiber square]
\label{prop:r5-incidence}
The curve $X_f$ is geometrically integral of bidegree $(1,3)$ and
arithmetic genus zero.  Projection to the root coordinate is
birational, while projection to the pencil coordinate gives a
degree-three morphism
\[
       \phi_f:\PP^1\longrightarrow\PP^1.
\]
If $\phi_f$ is separable, division of
\[
 c_1(t_1)c_2(t_2)-c_2(t_1)c_1(t_2)
\]
by the diagonal bracket defines a residual curve
$Y_f\subset\PP^1\times\PP^1$ of bidegree $(2,2)$ and arithmetic genus
one.  Its geometric components are the geometric-monodromy orbits on
ordered pairs of distinct roots.
\end{proposition}

\begin{proof}
A component of $X_f$ of bidegree $(0,e)$ would be a common factor of
$c_1,c_2$; hence none exists.  The bidegree genus formula gives genus
zero.  A root determines the unique pencil member containing it, so the
root projection is birational.  For a separable map the diagonal factor
in the fiber-square determinant has multiplicity one.  Removing its
bidegree $(1,1)$ from the fiber square leaves bidegree $(2,2)$, whose
arithmetic genus is one.  The standard fiber-square/monodromy
correspondence identifies its components with orbits on distinct
ordered sheets.
\end{proof}

\begin{lemma}[Cyclic stratum]\label{lem:cyclic}
Suppose $\phi_f$ is geometrically cyclic.
\begin{enumerate}[label=\textup{(\roman*)}]
\item In characteristic different from three there are two ramification
      points.  If they are rational, the corresponding osculating-plane
      intersection is deep exactly when $q\equiv2\pmod3$; if they are
      conjugate, it is deep exactly when $q\equiv1\pmod3$.
\item In characteristic three all osculating planes meet in
      $n=(0:0:1:0:0)$, whose pencil is
      $\langle T^3,U^3\rangle$; this fixed point is deep.
\item The remaining wild characteristic-three pencils have the form
      $\langle U^3,T^3+aTU^2\rangle$.  They are deep exactly when $-a$
      is nonsquare, forming one orbit of size $(q^2-1)/2$.
\end{enumerate}
\end{lemma}

\begin{proof}
\emph{Tame characteristic.}
Assume that the characteristic is not three.  Riemann--Hurwitz
gives two totally ramified points $t_1,t_2$.  The confluent Hankel
criterion places $f$ in
$\operatorname{Osc}(t_1)\cap\operatorname{Osc}(t_2)$; after transporting
the pair to $\{0,\infty\}$, the pencil is
$\langle T^3,U^3\rangle$ and a geometric deck generator is
$t\mapsto\zeta t$.

For a rational ramification pair, Frobenius fixes the points and the
deck map descends exactly when $\zeta^q=\zeta$, namely
$q\equiv1\pmod3$.  Thus its complement is deep exactly for
$q\equiv2\pmod3$.  The setwise stabilizer of the ordered normal form is
the split dihedral group of order $2(q-1)$, so the orbit size is
\[
       \frac{q(q^2-1)}{2(q-1)}=\frac{q(q+1)}2.
\]
For a conjugate pair Frobenius interchanges the fixed points and
inverts the multiplier; descent is instead $q\equiv2\pmod3$.  The
nonsplit dihedral stabilizer has order $2(q+1)$ and the orbit size is
$q(q-1)/2$.

\emph{The fixed wild point.}
Suppose the characteristic is three.  The second Hasse derivative
of
\[
       \nu(t)=(1,t,t^2,t^3,t^4)
\]
is $(0,0,1,3t,6t^2)=e_2$, including in the reversed infinity chart.
Thus every osculating plane contains $n=[e_2]$.  Substitution shows that
$n$ is $\PGL_2(q)$-fixed and
$W_n=\langle T^3,U^3\rangle$, so it is deep.

\emph{The remaining wild orbit.}
A remaining separable cyclic cover has one wild totally ramified
point.  Sending it to infinity and comparing coefficients under a
nontrivial translation gives
\[
 W_f=\langle U^3,T^3+\alpha T^2U+\beta TU^2\rangle,\qquad
 \alpha=0,\quad \kappa^2=-\beta.
\]
It is therefore represented by
\[
 f_a=e_2-ae_4,\qquad
 W_{f_a}=\langle U^3,T^3+aTU^2\rangle.
\]
The affine members are $z^3+az+s$.  Such a member has three rational
roots exactly when the additive map $z\mapsto z^3+az$ has a nonzero
rational kernel, equivalently when $-a$ is a square.  Hence the
nonsquare parameters form the deep orbit.  Its stabilizer is
$t\mapsto\pm t+\beta$, of order $2q$, and orbit--stabilizer gives
$(q^2-1)/2$.
\end{proof}

The count of split members of the pencil is not merely positive or
zero: it is determined exactly by the fibre square, and the following
identity is what drives both the threshold and the shape of the
exceptions.

\begin{proposition}[Exact split-witness count]\label{prop:r5-count}
Let $\deg\gcd(W_f)=0$ and let $\phi_f$ be separable; the characteristic
is arbitrary.
Write
\begin{itemize}
\item $N_f$ for the number of completely split squarefree members of
      $W_f$,
\item $d_2$ for the number of members with a rational double root and a
      distinct rational simple root,
\item $d_3$ for the number of members that are perfect cubes of a
      rational linear form.
\end{itemize}
Then
\begin{equation}\label{eq:r5-count}
       \#Y_f(\F_q)=6N_f+3d_2+d_3 .
\end{equation}
In particular $f$ is split-free if and only if
$\#Y_f(\F_q)=3d_2+d_3$.
\end{proposition}

\begin{proof}
Each $x\in\PP^1(\F_q)$ lies on exactly one member $g_x$ of the pencil,
because the pencil is basepoint-free; write $g_x=\ell_xh_x$ with
$\ell_x$ the linear form at $x$.  The fibre of $Y_f$ over $x$ under the
first projection is the set of rational roots of $h_x$, so
$\#Y_f(\F_q)=\sum_x\#\{\text{rational roots of }h_x\}$, and it suffices
to total that sum member by member.  A member with three distinct
rational roots is met at each of them, and there $h_x$ has the two
remaining roots, contributing $3\cdot2=6$.  A member with a rational
double root $r$ and distinct rational simple root $s$ is met at both:
over $r$ the residual is $\ell_r\ell_s$ with two distinct rational
roots, over $s$ it is $\ell_r^2$ with one, contributing $2+1=3$.  A
perfect cube is met only at its root, where the residual is a square,
contributing $1$.  A member with one rational root and an irreducible
quadratic cofactor contributes $0$ there, and a member with no rational
root is met at no rational $x$.  Summing gives \eqref{eq:r5-count}.
\end{proof}

The proof exhibits the fibre square as a double cover, and that is what
identifies it with the genus-one curve of the incidence literature.

\begin{proposition}[The fibre square is the incidence curve]
\label{prop:r5-fibre-is-elliptic}
Keep the hypotheses of Proposition~\ref{prop:r5-count} and let
$D_f(x)=\operatorname{disc}(h_x)/4$.  The first projection realizes
$Y_f$ as the double cover
\[
       w^2=D_f(x)
\]
of $\PP^1$, branched where $D_f$ vanishes.  If moreover $W_f$ misses the
twisted cubic and $\operatorname{char}\F_q\neq2,3$, then $D_f$ is the
discriminant quartic that Kaipa and Pradhan attach to the line, so the
normalization of $Y_f$ is their elliptic curve $E_{W_f}$ and
\[
       \#Y_f(\F_q)=\#E_{W_f}(\F_q).
\]
\end{proposition}

\begin{proof}
The fibre of $Y_f$ over $x$ under the first projection is the root set
of the residual quadratic $h_x$, and its two roots are interchanged by
the choice of a square root of $\operatorname{disc}(h_x)$; so the
projection is the stated double cover, and it is branched exactly where
that discriminant vanishes.  Kaipa and Pradhan form the discriminant of
the same residual quadratic and define $E_L$ as the non-singular model
of $w^2=D_L$ \cite[\S4--\S5]{KaipaPradhan2025}, taking as we do one quarter
of the discriminant of the same residual quadratic.  They therefore differ by
the square of the nonzero rational function $\lambda\in\F_q(x)^\times$
that relates the two normalizations of the residual quadratic,
\[
       D_f=\lambda^2D_L\quad\text{in }\F_q(x)^\times,
\]
with equality when the representative of $g_x$ is normalized as they
normalize it.  Equality of square classes in the function field, rather
than pointwise at rational $x$, is what rules out a quadratic twist and
leaves the cover unchanged.  For the point counts,
note that on that stratum $D_f$ has nonzero discriminant, so the double
cover displayed above is already smooth and is its own normalization,
and $Y_f$ is smooth there as well; the two smooth projective models are
then isomorphic and their counts agree.  Away from that stratum the
identification is birational and the counts may differ at a rational
singularity, which is why the exact count of
Proposition~\ref{prop:r5-count} is stated for $\#Y_f(\F_q)$ rather than
for the count on a normalization.
\end{proof}

\begin{lemma}[Branch budget]\label{lem:r5-branch}
Under the hypotheses of Proposition~\ref{prop:r5-count},
\[
       d_2+2d_3\leq4,
\]
so a split-free $f$ has $\#Y_f(\F_q)\leq12$, with equality only when
$d_2=4$ and $d_3=0$.
\end{lemma}

\begin{proof}
Riemann--Hurwitz gives different degree four for the separable
degree-three map $\phi_f$.  A member with a double root is a simple
branch value; it spends different degree one when the ramification is
tame and at least two when it is wild, which happens exactly in
characteristic two.  A member that is a perfect cube is totally
ramified; it spends two when tame and at least three when wild, which
happens exactly in characteristic three.  Distinct members are distinct
branch values, so $d_2+2d_3\leq4$ in every characteristic, and
characteristic two gives the sharper $d_2+d_3\leq2$.  Substituting into
\(\#Y_f(\F_q)=3d_2+d_3\) and maximizing over the integer points of
that region gives $12$, attained only at $(d_2,d_3)=(4,0)$.
\end{proof}

\begin{corollary}[Split members of an $S_3$ pencil]\label{cor:r5-equidistribution}
If $\phi_f$ has geometric monodromy $S_3$, then
\[
 \frac{q+1-2\sqrt q-12}{6}\;\leq\;N_f\;\leq\;\frac{q+1+2\sqrt q}{6}.
\]
\end{corollary}

\begin{proof}
Proposition~\ref{prop:r5-count} gives $6N_f=\#Y_f(\F_q)-3d_2-d_3$, and
Lemma~\ref{lem:r5-branch} bounds $3d_2+d_3$ between $0$ and $12$.  The
$S_3$ hypothesis makes $Y_f$ geometrically integral of arithmetic genus
one, so Aubry--Perret confines $\#Y_f(\F_q)$ to
$q+1\pm2\sqrt q$ \cite[p.~468]{AubryPerret1995}.
\end{proof}

The main term $(q+1)/6$ is the $S_3$ splitting density.  A completely
split fibre is one whose Frobenius class is the identity, of density
$1/|S_3|=1/6$ among the $q+1$ fibres, in the sense of the
factorization-type correspondence of \cite[Thm.~3.6]{Wang2026}.
Corollary~\ref{cor:r5-equidistribution} makes that quantitative: split
members occur with the expected density $1/6$ up to a Weil-scale error
and a bounded branch correction.  Split-freeness is the extreme
small-field case in which the lower bound has turned nonpositive, which
is why the exceptional fields are small ones.  Lemma~\ref{lem:s3}
records only the sign of this statement.  We claim the density only for
the completely split class; the other factorization types are not
counted here.

\begin{lemma}[$S_3$ stratum]\label{lem:s3}
If $\phi_f$ has geometric monodromy $S_3$ and $q\geq20$, then the pencil
contains a completely split squarefree cubic.  For odd $q$ the same
conclusion holds for every $q\geq20$ without the monodromy hypothesis
whenever $Y_f$ is geometrically integral.
\end{lemma}

\begin{proof}
By Proposition~\ref{prop:r5-incidence}, components of $Y_f$ correspond
to monodromy orbits on distinct ordered roots.  The $S_3$ action is
two-transitive, so $Y_f$ is geometrically integral and has arithmetic
genus one.  The Aubry--Perret singular-curve bound, in the form
\[
 \bigl|\#C(\F_q)-(q+1)\bigr|\leq 2\pi_C\sqrt q
\]
for a reduced geometrically integral projective curve of arithmetic
genus \(\pi_C\), gives
\[
       \#Y_f(\F_q)\geq q+1-2\sqrt q
\]
\cite[p.~468]{AubryPerret1995}.  Suppose $f$ were split-free.  Then
$\#Y_f(\F_q)\leq12$ by Lemma~\ref{lem:r5-branch}, so
$q+1-2\sqrt q\leq12$, that is $(\sqrt q-1)^2\leq12$ and
$q\leq(1+2\sqrt3)^2<20$.
\end{proof}

\begin{corollary}[Forced fibre-square invariants]\label{cor:r5-forced}
Let $q$ be odd with $17\leq q\leq19$ and let $f$ be split-free with
$Y_f$ geometrically integral.  Then
\[
 \#Y_f(\F_q)=12,\qquad d_2=4,\qquad d_3=0 .
\]
\end{corollary}

\begin{proof}
Lemma~\ref{lem:r5-branch} restricts $\#Y_f(\F_q)=3d_2+d_3$ to the
values $0,1,2,3,4,6,7,9,12$.  Aubry--Perret forces
$\#Y_f(\F_q)\geq q+1-2\sqrt q$, which exceeds $9$ for $q\geq17$; the
only admissible value left is $12$, and by Lemma~\ref{lem:r5-branch}
that value is attained only at $(d_2,d_3)=(4,0)$.
\end{proof}

Proposition~\ref{prop:r5-count} is a characteristic-uniform counting
identity, and in characteristic other than two and three its generic
specialization is an incidence count in the twisted-cubic line
literature.  When $W_f$
misses the twisted cubic, so that $d_3=0$, and
$\operatorname{char}\F_q\neq2,3$, the quantity $d_2$ is the invariant
$\eta_L$ attached to the binary quartic of the line and $\#Y_f(\F_q)$
is the point count of the associated elliptic curve, and
\eqref{eq:r5-count} becomes
\[
       N_f=\frac{\#E_L(\F_q)-3\eta_L}{6},
\]
which is the incidence count of \cite{KaipaPradhan2025}.  We cite it in
the form derived in their proof, namely
$|S\cap\mathcal O_3|=(\nu_L-\eta_L)/3$ together with
$\nu_L=(\#E_L(\F_q)-\eta_L)/2$
\cite[Prop.~4.5(3) and Thm.~5.1]{KaipaPradhan2025}; the denominator
printed in the displayed statement of their Theorem~1.3(3) is $3$, which
is a typographical slip, since $6$ is what their proof gives and what
makes their five incidence counts sum to $q+1$.  The statement above drops
the restriction $\operatorname{char}\F_q\neq2,3$ altogether, brings in
lines meeting the curve through the $d_3$ term, and needs no invariant
theory of quartics: its proof counts rational roots in fibres and uses
neither a discriminant nor a division by two.  What does stay inside
$\operatorname{char}\F_q\neq2,3$ is the comparison itself, through
Proposition~\ref{prop:r5-fibre-is-elliptic}, since the elliptic model is
built from a discriminant.  An
exhaustive check of \eqref{eq:r5-count} and of
Lemma~\ref{lem:r5-branch} over every point of $\mathrm{PG}(4,q)$ for the
prime powers $4\leq q\leq32$, in every characteristic, is recorded in the
supplement.

\begin{proposition}[Exhaustion and finite bridge]\label{prop:r5-bridge}
The preceding cases exhaust every redundancy-five pencil for
$q\geq20$.  Certificate R5 closes exactly
\[
       q\in\{7,8,9,11,13,16,17,19\}.
\]
It finds sporadic $S_3$-stratum orbits exactly at
$q=7,8,9,11,13,17,19$ and none at $q=16$.
\end{proposition}

\begin{proof}
The gcd degree is zero, one, or two.  Propositions
\ref{prop:r5-gcd2} and \ref{prop:r5-gcd1} settle the positive-gcd
cases.  In the trivial-gcd case the cubic map is inseparable or
separable.  The inseparable case occurs only in characteristic three
and the overlap equations give the all-cube nucleus pencil.  A
separable transitive cubic has geometric monodromy $C_3$ or $S_3$,
settled by Lemmas~\ref{lem:cyclic} and \ref{lem:s3}.  The stated finite
residue, canonical representatives, stabilizers, histograms, and
Frobenius links are exactly the domain of Certificate R5.
\end{proof}

\begin{remark}[Relation to the twisted-cubic line literature]
\label{rem:r5-line-literature}
Two conventions for ``the line of $\mathrm{PG}(3,q)$ attached to a pencil
of binary cubics'' are in use, and they are exchanged by a polarity, so
the class names below must be read in the right one.  G\"unay and
Lavrauw, and Blokhuis, Pellikaan, and Sz\H onyi, send a cubic to a
\emph{plane}: the twisted cubic is a curve of points, a cubic is the
plane meeting it in that cubic's divisor, and a pencil of cubics becomes
a pencil of planes, whose axis is the line
\cite[\S2]{GunayLavrauw2022}, \cite[Prop.~5.5]{BPS2022}.  Kaipa and
Pradhan send a cubic to a \emph{point}, so that the twisted cubic
consists of the perfect cubes and the pencil \emph{is} the line
\cite[\S1]{KaipaPradhan2025}.  The bridge is the symplectic polarity
$\sigma$ of $\mathrm{PG}(3,q)$, which carries each point of the twisted
cubic to its osculating plane and carries the chords of the twisted
cubic to the axes of the osculating developable
\cite[\S2]{GunayLavrauw2022}.

For us the translation is concrete.  The osculating plane at the point
$\lambda^3$ of the twisted cubic is $\lambda\Sym^2E^\vee$, so
\[
 \operatorname{Osc}(\lambda_1)\cap\operatorname{Osc}(\lambda_2)
       =\lambda_1\lambda_2\langle T,U\rangle
\]
is exactly a quadratic-gcd pencil, which is therefore an axis in the
point convention; by \cite[Lem.~2]{GunayLavrauw2022} its polar is the
chord joining $\lambda_1^3$ and $\lambda_2^3$, which is what the plane
convention calls it.  Symmetrically $\langle T^3,U^3\rangle$ joins two
points of the twisted cubic and so is a chord in the point convention
and an axis in the plane convention.  The class names in this remark
follow the plane convention of the source being cited; the split members
of a pencil are of course the same either way.

Under \eqref{eq:r5-orbit-criterion} every stratum above is a class of
lines of $\mathrm{PG}(3,q)$ in the sense of Blokhuis, Pellikaan, and
Sz\H onyi, who partition the lines into classes $O_1,\ldots,O_8$ and
determine for $q\geq23$, in every characteristic, which classes lie in
a plane meeting the twisted cubic in three rational points
\cite[Thm.~7.1 and Prop.~7.4]{BPS2022}.  The correspondence is exact.
Our quadratic-gcd trichotomy of Proposition~\ref{prop:r5-gcd2} is their
degree-one case: real chords $O_1$ are shallow, tangents $O_2$ give the
tangent family, imaginary chords $O_3$ give the conjugate-secant
family.  The linear-gcd unisecants of
Proposition~\ref{prop:r5-gcd1} are their $O_4$ and $O_5$.  The two
cyclic cases of Lemma~\ref{lem:cyclic}(i) are their real axes $O_1'$
and imaginary axes $O_3'$, with the same congruences mod three; the
characteristic-three nucleus of Lemma~\ref{lem:cyclic}(ii) is their
axis $O_7(3)$; the wild pencils of Lemma~\ref{lem:cyclic}(iii) are
their $O_{8.1}(3)^{\pm}$ and $O_{8.2}(3)$.  Lemma~\ref{lem:s3} is their
$O_6$, and their Remark 6.12 reaches the threshold $q\geq23$ by the
same route: the double point scheme of the degree-three map is a
genus-one curve, at most twelve points must be excluded, and the
Hasse--Weil bound then produces a fibre with three distinct rational
points.  Two differences remain.  Aubry--Perret in place of
Hasse--Weil drops their simplicity hypothesis, which is what lets one
estimate carry the singular and wild cases uniformly; and the exact
count of Proposition~\ref{prop:r5-count} replaces the deletion
bookkeeping by an identity, which lowers the threshold to $q\geq20$
and, through Corollary~\ref{cor:r5-forced}, pins the fibre square of
every remaining exception at $q=17,19$.  Independent confirmation of the same table comes from the
classification of degree-three permutation rational functions
\cite{FerragutiMicheli2020}, as recorded in \cite[Rem.~7.6]{BPS2022}.

Exact counts, rather than the vanishing test, are available too.  For
the ten non-generic line classes Davydov, Marcugini, and Pambianco give
the exact number of planes of each orbit through each line for
$q\geq5$ \cite[Thm.~3.3]{DMP2022}, building on their line-orbit
classification \cite{DMP2023,DMP2021}; G\"unay and Lavrauw give the
point-orbit and plane-orbit distributions of the same classes for odd
$q$ prime to three \cite{GunayLavrauw2022}.  For the generic class and
$\operatorname{char}\F_q\neq2,3$, Kaipa and Pradhan attach to a line
$L$ a binary quartic $\varphi_L$ and an elliptic curve $E_L$, and count
$|W_f\cap\mathcal O_3|$ in terms of $\#E_L(\F_q)$ and the invariant
$\eta_L\in\{0,1,2,4\}$ recording the rational factorization type of
$\varphi_L$.  That count is the specialization of
Proposition~\ref{prop:r5-count} discussed after
Corollary~\ref{cor:r5-forced}, and
Proposition~\ref{prop:r5-fibre-is-elliptic} identifies their curve with
our fibre square.  By \eqref{eq:r5-orbit-criterion} a generic line is
split-free exactly when $\#E_L(\F_q)=3\eta_L\leq12$, which the Hasse
bound $\#E_L(\F_q)\geq q+1-2\sqrt q$ forbids for $q\geq20$ and permits
only for $q\leq19$.  This is not a second route to
Lemma~\ref{lem:s3} but the same one in other coordinates:
Proposition~\ref{prop:r5-fibre-is-elliptic} identifies $E_L$ with $Y_f$,
so both arguments apply a Weil-type bound to a single genus-one curve.
What the elliptic form adds is not independence but an explanation, since
$3\eta_L\leq12$ is visibly a bounded quantity while $\#E_L(\F_q)$ grows
with $q$.  The generic class is settled in
characteristic three in \cite{KaipaPradhan2025char3} and in
characteristic two by Ceria and Pavese \cite{CeriaPavese2023}.

What this section contributes is therefore not the pencil-level
classification but its syndrome-level source and consequence.  The
passage $f\mapsto W_f$ of Lemma~\ref{lem:hankel} is not present in that
literature, and it is not surjective onto the lines: the
characteristic-two inseparable unisecants, split-free as lines, are
realized by no rank-two syndrome, as the proof of
Proposition~\ref{prop:r5-gcd1} shows.  Neither do incidence
distributions produce syndromes; Table~\ref{tab:r5sporadic} records
canonical representatives in divided-power syndrome coordinates with
their orbit sizes, stabilizers, and Frobenius fusion, and extracting
such a list from \cite[Prop.~4.5 and Thm.~5.1]{KaipaPradhan2025} would still require
evaluating $\#E_L$ and $\eta_L$ over all $2q-3+\mu$ generic line orbits.
The promotion from split-free to deep, through
Proposition~\ref{prop:r5-radius}, is likewise outside that literature,
which contains no covering-radius or deep-hole statement.
\end{remark}

\begin{longtable}{c>{\ttfamily\small}lrrl}
\caption{Complete sporadic orbit inventory at redundancy five.  A
representative $[a_0,\ldots,a_4]$ is in the divided-power syndrome
coordinates of Section~\ref{sec:dictionary}.  The last column records
the nontrivial Frobenius fusion; a dash means the orbit is fixed.}
\label{tab:r5sporadic}\\
\toprule
$q$ & \normalfont representative & orbit size & stabilizer & Frobenius\\
\midrule
\endfirsthead
\toprule
$q$ & \normalfont representative & orbit size & stabilizer & Frobenius\\
\midrule
\endhead
\bottomrule
\endfoot
7  & [0,1,0,0,3] & 28  & 12 & --\\
7  & [0,1,0,0,2] & 112 & 3  & --\\
7  & [0,0,1,0,3] & 168 & 2  & --\\
7  & [0,1,0,3,2] & 168 & 2  & --\\
7  & [1,0,0,2,1] & 168 & 2  & --\\
8  & [0,1,1,0,2] & 252 & 2  & $\to[0,1,1,0,4]$\\
8  & [0,1,1,0,4] & 252 & 2  & $\to[0,1,1,0,6]$\\
8  & [0,1,1,0,6] & 252 & 2  & $\to[0,1,1,0,2]$\\
9  & [1,0,0,1,2] & 180 & 4  & --\\
9  & [0,1,0,4,1] & 360 & 2  & $\leftrightarrow[0,1,0,4,4]$\\
9  & [0,1,0,4,4] & 360 & 2  & $\leftrightarrow[0,1,0,4,1]$\\
11 & [1,0,0,1,10]& 330 & 4  & --\\
11 & [1,0,0,1,1] & 660 & 2  & --\\
13 & [0,1,0,0,4] & 182 & 12 & --\\
13 & [1,0,0,4,7] & 546 & 4  & --\\
17 & [1,0,0,1,5] & 1224& 4  & --\\
19 & [0,1,0,0,4] & 570 & 12 & --\\
\end{longtable}

The certificate additionally records the complete cubic-pencil member
histogram of each representative.  Repeated orbit sizes in
Table~\ref{tab:r5sporadic} are therefore visibly distinguished without
requiring a working-directory file.

\begin{table}[H]
\centering
\small
\caption{Independent R5 completeness totals.  The first equality is
the direct syndrome enumeration versus the orbit--stabilizer sum; the
second decomposes the same total into the formula in
Theorem~\ref{thm:r5} and the sporadic rows above.}
\label{tab:r5-sanity}
\begin{tabular}{crrr}
\toprule
$q$ & direct $=$ orbit sum & nonsporadic & sporadic\\
\midrule
7  & $889=889$   & 245  & 644\\
8  & $1116=1116$ & 360  & 756\\
9  & $1391=1391$ & 491  & 900\\
11 & $1848=1848$ & 858  & 990\\
13 & $2080=2080$ & 1352 & 728\\
16 & $2432=2432$ & 2432 & 0\\
17 & $4131=4131$ & 2907 & 1224\\
19 & $4541=4541$ & 3971 & 570\\
\bottomrule
\end{tabular}
\end{table}

\begin{remark}[No one-column MDS extension at this radius]
\label{rem:q8-no-extension}
At \(q=8\) the split-free count of Table~\ref{tab:r5-sanity} is
\(1116\), splitting as \(360\) nonsporadic and \(756\) sporadic
directions.  These are deep holes of \(\PRS_{\F_8}(4)\) by
Proposition~\ref{prop:r5-radius} and
Corollary~\ref{cor:splitfree}; they are not one-column MDS extensions,
and no such extension exists.

The two notions differ by one column.  A syndrome direction is a deep
hole when it lies in the span of no \(r-2\) parity-check columns and
the radius is \(r-1\); adjoining a point to the \((q+1)\)-arc of the
normal rational curve keeps it an arc, equivalently extends the dual
code by one column, only when that point lies in the span of no
\(r-1\) columns.  At redundancy five the first condition gives distance
at least four and Proposition~\ref{prop:r5-radius} gives radius exactly
four, so no direction achieves the distance five that an extension
requires.  D\"ur's equivalence states this as a biconditional: the
covering radius is \(r-1\) precisely when the arc of every normal
rational curve in \(\mathrm{PG}(r-1,q)\) is complete, so there is no MDS
extension of the dual by one digit
\cite{Dur1994}, as stated by Kaipa \cite[Sec.~IV]{Kaipa2017}.  The same
reading applies at every redundancy classified here, since each radius
gate proves \(\rho=r-1\).

Independently, a \([10,5,6]_8\) MDS code would be a \(10\)-arc in
\(\mathrm{PG}(4,8)\), and the largest arc there has \(q+1=9\) points.
\end{remark}
